% Vietnamese translation for OI Public Library
% Source-File: IOI中国国家候选队论文/2009/徐源盛/对一类动态规划问题的研究.ppt
\documentclass[11pt]{article}
\usepackage{oipl-vn}

\newcommand{\Oh}{\mathcal{O}}

\title{Nghiên cứu về một lớp bài toán quy hoạch động\\{\large Ghi chú theo slide}}
\author{Từ Nguyên Thịnh\\Trường Trung học số 1 Trường Sa, Hồ Nam}
\date{2009}

\begin{document}
\maketitle

\origtitle{A study of a class of dynamic programming problems}
\sourcepath{IOI-China-Candidate-Team-Papers/2009/Xu-Yuansheng/future-cost-dp-slides.ppt}

\section{Chủ đề}

Slide bàn về những bài quy hoạch động mà quyết định hiện tại tạo ra chi phí
hoặc ảnh hưởng trong tương lai. Thay vì chỉ ghi nhận trạng thái hiện tại, ta
phải mô hình hóa phần ảnh hưởng chưa trả hết để công thức chuyển giữ được tính
đúng đắn.

\section{Hai kiểu ảnh hưởng}

\begin{enumerate}
  \item Ảnh hưởng của quyết định hiện tại chỉ đi qua chính đối tượng vừa chọn.
  Khi đó trạng thái có thể gọn, thường chỉ cần vị trí cuối hoặc cấu hình biên.
  \item Ảnh hưởng còn phụ thuộc vào tình huống tương lai. Khi đó phải thêm
  thông tin phụ để mô tả khoản chi phí treo, hoặc đổi thứ tự xử lý để biến
  tương lai thành quá khứ.
\end{enumerate}

\section{Các ví dụ định hướng}

\textit{Những quả cầu của Sue} và \textit{Bài toán sắp xếp thứ hạng} minh họa
cách giữ đủ thông tin để quyết định sau này tính đúng chi phí. \textit{Xóa
khối} và \textit{Logistics Olympic} cho thấy khi trạng thái nhìn thấy ban đầu
không đủ, việc thêm tham số hoặc đổi mô hình có thể làm công thức chuyển rõ
ràng hơn.

\section{Ghi chú}

Khi đọc slide, nên theo dõi bất biến của trạng thái: nó đã đại diện cho toàn
bộ lịch sử cần thiết chưa, và mọi chi phí tương lai có thể được tính chỉ từ
trạng thái đó hay không. Đây là tiêu chuẩn quan trọng hơn hình thức công thức
truy hồi.

\end{document}
